% GENERATED FILE. Edit canonical lessons, then run npm run generate:books.
% canonical-source-hash: fnv1a64:3e27aa41d67c3efd
% canonical-lessons: LA-C53-iterum, LA-C53-magister, LA-C53-discipulus, LA-C53-schola, LA-C53-liber

\chapter{Iterum, Magister, Discipulus, Schola, Liber}
\label{ch:la-schola}
\hlchaptermodality{53}

\begin{chapteropening}
\textbf{By the end of this chapter:} \emph{I can ask for something to be said again, and name the teacher, the pupil, the school and the book.}

The fifth word of the group, taken apart against the four before it.
\end{chapteropening}

\section[iterum]{iterum — and the repair move this book could not make}

\label{lesson:LA-C53-iterum}



\begin{quote}
Six recalls before the new word, at growing distances. They are not decoration: each one is the retrieval window this book measures, and each names a word far enough back that saying it is work.

\begin{itemize}
\raggedright
  \item \emph{Recall:} say \emph{tuus} — \textbf{R1}, one lesson back, before it decays
  \item \emph{Recall:} say \emph{ubi} — \textbf{R2}, the first real retrieval, five to fifteen lessons back
  \item \emph{Recall:} say \emph{possum} — \textbf{R3}, durable, twenty to sixty lessons back
  \item \emph{Recall:} say \emph{audiō} — \textbf{R3}, durable again, a second word from the same distance
  \item \emph{Recall:} say \emph{propediem tē vidēbō} — \textbf{R4}, recognition at distance, eighty lessons back or more
  \item \emph{Recall:} say \emph{crās tē vidēbō} — \textbf{R4}, recognition at distance, a second word from that far back
\end{itemize}
\end{quote}

\begin{cousinweb}
\textbf{iterum} — \textquotedblleft{}\textbf{again}.\textquotedblright{}

This one word closes the largest hole in this book. The \emph{quaesō} lesson gave you \textquotedblleft{}please\textquotedblright{}; the eight-verb chapter gave you \textbf{dīcō}, \textquotedblleft{}I say.\textquotedblright{} Fifty-two chapters went by without a way of joining them, and here it is:

\begin{quote}
\textbf{Iterum, quaesō.} — Again, please. \textbf{Quid dīcis?} — What are you saying? \textbf{Lentē, quaesō.} — Slowly, please.
\end{quote}

That is a \textbf{repair}: the move you make when you did not catch what somebody said. Every language teaches it in its first week, and this book — which could greet, thank, apologise and say please within four chapters — could not until now ask anyone to repeat anything.

\emph{Iterum} is built on the same \emph{\textbf{*i-}} stem as \textbf{is} from the pointing chapter, with an old comparative suffix: literally \textquotedblleft{}the other of the two,\textquotedblright{} hence \textquotedblleft{}a second time.\textquotedblright{} English \textbf{iterate} and \textbf{reiteration} are loans from it.
\end{cousinweb}

\subsection*{Your turn}
Say these aloud:

\begin{itemize}
\raggedright
  \item \emph{iterum}
  \item every word this chapter has given you so far, in order
\end{itemize}

\begin{tcolorbox}[breakable,colback=teal!4,colframe=teal!35!black,title={Before you move on}]
How do you ask a Roman to say it again? (\emph{\textbf{Iterum, quaesō}}.) Which two earlier lessons did the parts come from? (\textbf{The \emph{quaesō} lesson}, and the \textbf{eight-verb chapter} for \emph{dīcō}.)
\end{tcolorbox}
\section[magister, magistrī]{magister — the greater one, against minister, the lesser}

\label{lesson:LA-C53-magister}



\begin{quote}
Six recalls before the new word, at growing distances. They are not decoration: each one is the retrieval window this book measures, and each names a word far enough back that saying it is work.

\begin{itemize}
\raggedright
  \item \emph{Recall:} say \emph{iterum} — \textbf{R1}, one lesson back, before it decays
  \item \emph{Recall:} say \emph{cūr} — \textbf{R2}, the first real retrieval, five to fifteen lessons back
  \item \emph{Recall:} say \emph{volō} — \textbf{R3}, durable, twenty to sixty lessons back
  \item \emph{Recall:} say \emph{dormiō} — \textbf{R3}, durable again, a second word from the same distance
  \item \emph{Recall:} say \emph{propediem tē vidēbō} — \textbf{R4}, recognition at distance, eighty lessons back or more
  \item \emph{Recall:} say \emph{crās tē vidēbō} — \textbf{R4}, recognition at distance, a second word from that far back
\end{itemize}
\end{quote}

\begin{cousinweb}
\textbf{magister} — \textquotedblleft{}\textbf{a teacher, a master}.\textquotedblright{}

The word is one half of a matched pair and the other half is \textbf{minister}. \emph{Magis} means \textquotedblleft{}more,\textquotedblright{} \emph{minus} means \textquotedblleft{}less,\textquotedblright{} and both take the same old contrastive suffix \emph{-ter}: the \textbf{greater} one and the \textbf{lesser} one. A Roman \emph{magister} outranks a Roman \emph{minister}, and a modern government minister is still, etymologically, a servant.

The English family is enormous, and all of it is one Latin word:

\textbf{master}, \textbf{mister}, \textbf{Mr}, \textbf{mistress}, \textbf{Mrs}, \textbf{Miss}, \textbf{magistrate}, \textbf{maestro}, \textbf{magisterial} — and Spanish \emph{maestro}, French \emph{maître}, Italian \emph{maestro}.

And \emph{magis} itself you have met: The \emph{sed} lesson recorded that \emph{sed} died and \emph{magis} became Spanish \emph{mas} and Italian \emph{ma}, \textquotedblleft{}but.\textquotedblright{} The same \textquotedblleft{}more\textquotedblright{} is inside both.
\end{cousinweb}

\subsection*{Your turn}
Say these aloud:

\begin{itemize}
\raggedright
  \item \emph{magister, magistrī}
  \item every word this chapter has given you so far, in order
\end{itemize}

\begin{tcolorbox}[breakable,colback=teal!4,colframe=teal!35!black,title={Before you move on}]
What is the pair \emph{magister} belongs to? (\emph{\textbf{Minister}} — greater and lesser.) Which earlier word is \emph{magis} also behind? (\textbf{The Romance \emph{but}} — Spanish \emph{mas}, Italian \emph{ma}.)
\end{tcolorbox}
\section[discipulus, discipulī]{discipulus — from learning, and the reason discipline means what it does}

\label{lesson:LA-C53-discipulus}



\begin{quote}
Six recalls before the new word, at growing distances. They are not decoration: each one is the retrieval window this book measures, and each names a word far enough back that saying it is work.

\begin{itemize}
\raggedright
  \item \emph{Recall:} say \emph{magister} — \textbf{R1}, one lesson back, before it decays
  \item \emph{Recall:} say \emph{quia} — \textbf{R2}, the first real retrieval, five to fifteen lessons back
  \item \emph{Recall:} say \emph{quī} — \textbf{R3}, durable, twenty to sixty lessons back
  \item \emph{Recall:} say \emph{sedeō} — \textbf{R3}, durable again, a second word from the same distance
  \item \emph{Recall:} say \emph{crās tē vidēbō} — \textbf{R4}, recognition at distance, eighty lessons back or more
  \item \emph{Recall:} say \emph{quōmodo} — \textbf{R4}, recognition at distance, a second word from that far back
\end{itemize}
\end{quote}

\begin{cousinweb}
\textbf{discipulus} — \textquotedblleft{}\textbf{a pupil}.\textquotedblright{}

Built on \textbf{discō}, \textquotedblleft{}I learn,\textquotedblright{} which is itself a reduplicated form of the same root as \emph{doceō}, \textquotedblleft{}I teach.\textquotedblright{} Learning and teaching are one word split in two, and the pupil is named for the learning half.

Then the word that grew out of it. \textbf{Disciplīna} originally meant no more than \textquotedblleft{}teaching, instruction\textquotedblright{} — what a \emph{discipulus} receives. Only later did it narrow to the training that keeps order, and later still to punishment. English \textbf{discipline} carries all three senses at once, and a \textbf{disciple} is the \emph{discipulus} of the Latin Bible.

Note the shape against the last lesson's. \textbf{Magister} and \textbf{discipulus} are the two ends of a Roman classroom, and Latin names one for rank and the other for what he does.
\end{cousinweb}

\subsection*{Your turn}
Say these aloud:

\begin{itemize}
\raggedright
  \item \emph{discipulus, discipulī}
  \item every word this chapter has given you so far, in order
\end{itemize}

\begin{tcolorbox}[breakable,colback=teal!4,colframe=teal!35!black,title={Before you move on}]
What verb is \emph{discipulus} built on? (\emph{\textbf{Discō}}, I learn.) What did \emph{disciplīna} first mean? (\textbf{Teaching} — not punishment.)
\end{tcolorbox}
\section[schola, scholae]{schola — a Greek word that first meant leisure}

\label{lesson:LA-C53-schola}



\begin{quote}
Six recalls before the new word, at growing distances. They are not decoration: each one is the retrieval window this book measures, and each names a word far enough back that saying it is work.

\begin{itemize}
\raggedright
  \item \emph{Recall:} say \emph{discipulus} — \textbf{R1}, one lesson back, before it decays
  \item \emph{Recall:} say \emph{meus} — \textbf{R2}, the first real retrieval, five to fifteen lessons back
  \item \emph{Recall:} say \emph{et} — \textbf{R3}, durable, twenty to sixty lessons back
  \item \emph{Recall:} say \emph{stō} — \textbf{R3}, durable again, a second word from the same distance
  \item \emph{Recall:} say \emph{quōmodo} — \textbf{R4}, recognition at distance, eighty lessons back or more
  \item \emph{Recall:} say \emph{bene} — \textbf{R4}, recognition at distance, a second word from that far back
\end{itemize}
\end{quote}

\begin{cousinweb}
\textbf{schola} — \textquotedblleft{}\textbf{a school}.\textquotedblright{}

Greek, and the meaning has turned completely over. \emph{Scholē} meant \textbf{leisure} — spare time, the opposite of work. Only a person with leisure could spend it on learned discussion, so the word moved to the discussion, then to the place the discussion happened, then to the building children are sent to.

A word that meant \emph{free time} now means \emph{the place you must be at nine}.

Latin took it as \textbf{schola}, and every daughter kept it: Spanish \emph{escuela}, Italian \emph{scuola}, French \emph{école}, Portuguese \emph{escola} — and each of them props up the \emph{sc-} with a vowel in front, exactly the way Portuguese props up \emph{escrever}. English \textbf{school}, \textbf{scholar} and \textbf{scholastic} are loans; \textbf{schedule} is not, and the resemblance is a coincidence.
\end{cousinweb}

\subsection*{Your turn}
Say these aloud:

\begin{itemize}
\raggedright
  \item \emph{schola, scholae}
  \item every word this chapter has given you so far, in order
\end{itemize}

\begin{tcolorbox}[breakable,colback=teal!4,colframe=teal!35!black,title={Before you move on}]
What did \emph{scholē} originally mean? (\textbf{Leisure}.) Why do Spanish and French put a vowel in front of \emph{sc-}? (\textbf{They prop up the cluster} — \emph{escuela}, \emph{école}.)
\end{tcolorbox}
\section[liber, librī]{liber — and the bark it was written on}

\label{lesson:LA-C53-liber}



\begin{quote}
Six recalls before the new word, at growing distances. They are not decoration: each one is the retrieval window this book measures, and each names a word far enough back that saying it is work.

\begin{itemize}
\raggedright
  \item \emph{Recall:} say \emph{schola} — \textbf{R1}, one lesson back, before it decays
  \item \emph{Recall:} say \emph{tuus} — \textbf{R2}, the first real retrieval, five to fifteen lessons back
  \item \emph{Recall:} say \emph{quandō} — \textbf{R3}, durable, twenty to sixty lessons back
  \item \emph{Recall:} say \emph{ambulō} — \textbf{R3}, durable again, a second word from the same distance
  \item \emph{Recall:} say \emph{quōmodo} — \textbf{R4}, recognition at distance, eighty lessons back or more
  \item \emph{Recall:} say \emph{bene} — \textbf{R4}, recognition at distance, a second word from that far back
\end{itemize}
\end{quote}

\begin{cousinweb}
\textbf{liber} — \textquotedblleft{}\textbf{a book}.\textquotedblright{}

The word for the thing this whole book is trying to let you read, and it arrives here, in the second-to-last chapter written so far.

\emph{Liber} first meant the \textbf{inner bark} of a tree — the surface Romans wrote on before papyrus was easy to get. The material became the text and the text became the book, the same journey English \emph{paper} made from \emph{papyrus} and \emph{volume} made from \emph{volvō}, \textquotedblleft{}I roll,\textquotedblright{} because a Roman book was a scroll.

Watch out for the homograph. \textbf{Līber} with a long \emph{ī} means \textquotedblleft{}\textbf{free},\textquotedblright{} and gives English \emph{liberty}, \emph{liberal} and \emph{deliver}. \textbf{Liber} with a short \emph{i} is the book. Two different words, four identical letters, and only the macron between them — the same lesson \textbf{hic} and \textbf{hīc} taught in the pointing chapter.

Spanish \emph{libro}, Italian \emph{libro}, French \emph{livre}, Portuguese \emph{livro}. English \textbf{library} and \textbf{libretto} (\textquotedblleft{}little book\textquotedblright{}) are loans.
\end{cousinweb}

\subsection*{Your turn}
Say these aloud:

\begin{itemize}
\raggedright
  \item \emph{liber, librī}
  \item every word this chapter has given you so far, in order
\end{itemize}

\begin{tcolorbox}[breakable,colback=teal!4,colframe=teal!35!black,title={Before you move on}]
What did \emph{liber} first mean? (\textbf{Tree bark} — the writing surface.) What separates \emph{liber} from \emph{līber}? (\textbf{The macron} — book against free.)
\end{tcolorbox}
